\section{Test-Run Audit Trail}\label{app:test}
This appendix details the safeguards of the one-time test evaluation summarised in Section~\ref{sec:test}
(script \texttt{scripts/run\_phase7.py}; post-hoc analysis of the saved outputs in
\texttt{scripts/phase7\_posthoc.py}; outputs in \texttt{results/phase7/run\_1/}).
\begin{itemize}\setlength\itemsep{0pt}
\item \emph{Frozen models.} The classifiers were refit on all development images with the hyper-parameters
selected most often in the nested cross-validation; the refit of the main forensic random forest is identical to
the deployed model's forest. The deployed calibrated models were used with their uncertain bands, and the deployed
localisers with their post-processing. Localisation baselines had their post-processing tuned on the validation
images only.
\item \emph{Dry runs.} Before the test run, the full script was exercised twice in a dry-run mode, with the
training images as ``development'' and the validation images as ``test''. It was then checked in an independent code review.
\item \emph{Preflight.} At run time a preflight confirmed that every model, feature file and cache matched the
current code and configuration.
\item \emph{Start record.} The start record in \texttt{results/phase7/test\_runs.log} holds fingerprints of 36
files (all project source files and \texttt{config.yaml}) and of 12 input files (the frozen split's hash file,
the models and the classification and localisation inputs they were built from).
\item \emph{Completion record.} Because of a path bug, run~1 wrote its completion record, with the hashes of
\texttt{summary.json} and \texttt{test\_classification.csv}, into \texttt{run\_1/test\_runs.log}. The record was
copied, unchanged and annotated, into the main log after the run, and the bug was then fixed. The fix changes the
script's fingerprint for any future run.
\item \emph{Output manifest.} A SHA-1 manifest of every output file of run~1 (\texttt{run\_1/MANIFEST.sha1}) was
recorded on the same day.
\item \emph{Post-run check.} A post-run check recomputed the code and input fingerprints and found nothing
changed since the run started.
\end{itemize}
The later uses of the test split, namely the post-hoc robustness, synthetic and external evaluations with the same
frozen models (Section~\ref{sec:robustness}) and the pre-registered, one-time comparison with an ELA-CNN baseline
(Section~\ref{sec:cnn}; new models, evaluated once), fed nothing back and are recorded in the same log
(Appendix~\ref{app:repro}).
